\section*{Appendix: Technical Details of the Bayesian Latent Factor Model}
\label{appendix:blf}

This appendix provides additional technical details on the Bayesian Latent Factor (BLF) model, including prior specification, inference algorithms, sensitivity analyses, and extended validation results.

\subsection*{A.1 Prior Specification and Justification}

\subsubsection*{A.1.1 Latent Factor Prior}

The latent factor $\theta_i$ represents the standardized composite score for model $i$:

\begin{equation}
\theta_i \sim \mathcal{N}(0, 1)
\end{equation}

This standard normal prior ensures identifiability: without it, the scale of $\theta$ would be confounded with the loadings $\lambda_b$. The prior is weakly informative---it centers the distribution but allows data to override if needed.

\subsubsection*{A.1.2 Benchmark-Specific Parameters}

\textbf{Intercepts ($\alpha_b$):} Account for benchmark difficulty and mean offset:
\begin{equation}
\alpha_b \sim \mathcal{N}(0, 2^2) \quad (\text{Variance} = 4)
\end{equation}

Since benchmarks are standardized (z-scores), we expect $\alpha_b \approx 0$, but allow deviations up to $\pm 2$ standard deviations to accommodate systematic biases.

\textbf{Loadings ($\lambda_b$):} Represent the importance/weight of each benchmark:
\begin{equation}
\lambda_b \sim \text{HalfNormal}(1)
\end{equation}

The half-normal prior enforces $\lambda_b > 0$ (benchmarks positively correlate with quality), reflecting the assumption of a positive manifold where all valid benchmarks measure the same underlying construct. It has mean $\sqrt{2/\pi} \approx 0.80$, which is reasonable for factor loadings. The prior is weakly informative: it allows loadings from 0.3 to 1.5 with high probability.

\textbf{Noise ($\sigma_b$):} Captures measurement error and benchmark-specific variance:
\begin{equation}
\sigma_b \sim \text{HalfNormal}(1)
\end{equation}

This prior allows for varying noise levels across benchmarks. Benchmarks with low $\sigma_b$ are more reliable and contribute more to the posterior.

\subsection*{A.2 Inference Algorithm}

We use the No-U-Turn Sampler (NUTS) \cite{hoffman2014nuts}, a variant of Hamiltonian Monte Carlo (HMC) that adaptively tunes the integration time. NUTS is the default sampler in PyMC and Stan due to its excellent performance on hierarchical models.

\subsubsection*{A.2.1 Sampler Configuration}

\begin{itemize}
    \item \textbf{Algorithm:} NUTS (No-U-Turn Sampler)
    \item \textbf{Chains:} 4 independent chains initialized from overdispersed starting values
    \item \textbf{Tuning:} 2,000 iterations (warmup) for mass matrix adaptation and step size tuning
    \item \textbf{Sampling:} 2,000 iterations per chain (8,000 total samples)
    \item \textbf{Target acceptance:} 0.97 (high to avoid divergences in tails)
    \item \textbf{Max tree depth:} 10 (default)
    \item \textbf{Initialization:} jitter=0.1 around MAP estimate
\end{itemize}

\subsubsection*{A.2.2 Convergence Diagnostics}

We assess convergence using multiple diagnostics:

\begin{enumerate}
    \item \textbf{Gelman-Rubin $\hat{R}$ statistic} \cite{gelman1992inference}: Compares within-chain and between-chain variance. We require $\hat{R} < 1.01$ for all parameters.
    
    \begin{equation}
    \hat{R} = \sqrt{\frac{\hat{V}}{W}}
    \end{equation}
    
    where $\hat{V}$ is the marginal posterior variance and $W$ is the within-chain variance.
    
    \item \textbf{Effective Sample Size (ESS)}: Accounts for autocorrelation in MCMC chains. We require ESS $> 400$ per chain (1,600 total) for reliable inference.
    
    \item \textbf{Trace plots}: Visual inspection of chain mixing and stationarity.
    
    \item \textbf{Divergent transitions}: We monitor for divergences, which indicate problematic posterior geometry. Target: 0 divergences.
    
    \item \textbf{Energy plot}: Bayesian Fraction of Missing Information (BFMI) should be $> 0.3$ \cite{betancourt2016diagnosing}.
\end{enumerate}

\subsection*{A.3 Model Validation}

\subsubsection*{A.3.1 Posterior Predictive Checks}

We validate model fit by comparing observed data to posterior predictive samples:

\begin{equation}
z_{i,b}^{\text{rep}} \sim \mathcal{N}(\alpha_b + \lambda_b \theta_i, \sigma_b^2)
\end{equation}

drawn from the posterior. Good fit requires:
\begin{itemize}
    \item $R^2 > 0.80$ between observed and predicted z-scores
    \item Residuals centered at 0 with no systematic patterns
    \item Coverage: 95\% of observed scores fall within 95\% predictive intervals
\end{itemize}

\subsubsection*{A.3.2 Cross-Validation}

We perform 5-fold cross-validation at the model level:
\begin{itemize}
    \item Hold out 20\% of models in each fold
    \item Fit BLF on remaining 80\%
    \item Predict held-out scores using posterior predictive distribution
    \item Metric: Mean Absolute Error (MAE) on held-out standardized scores
\end{itemize}

Results: MAE = 0.31 $\pm$ 0.04 (mean $\pm$ SE across folds), indicating good generalization.

\subsubsection*{A.3.3 External Validation}

We validate composite scores against external quality indicators:

\begin{table}[h]
\centering
\caption{External Validation of CCS}
\begin{tabular}{lcc}
\toprule
\textbf{External Metric} & \textbf{Spearman $\rho$} & \textbf{p-value} \\
\midrule
Chatbot Arena ELO (Coding) & 0.89 & $< 0.001$ \\
HuggingFace Model Downloads & 0.62 & $< 0.001$ \\
GitHub Stars (for open models) & 0.71 & $< 0.001$ \\
Human Preference (N=50 pairs) & 0.83 & $< 0.001$ \\
\bottomrule
\end{tabular}
\end{table}

\subsection*{A.4 Sensitivity Analyses}

\subsubsection*{A.4.1 Prior Sensitivity}

We test robustness to prior specification by varying prior hyperparameters:

\begin{table}[h]
\centering
\caption{Sensitivity to Prior Specification}
\begin{tabular}{lccc}
\toprule
\textbf{Prior Variant} & \textbf{Mean CCS} & \textbf{SD CCS} & \textbf{$\rho$ vs Default} \\
\midrule
Default & 77.3 & 12.1 & 1.00 \\
$\lambda \sim \text{HalfNormal}(0.5)$ & 77.1 & 12.3 & 0.998 \\
$\lambda \sim \text{HalfNormal}(2)$ & 77.4 & 12.0 & 0.999 \\
$\sigma \sim \text{HalfNormal}(0.5)$ & 77.5 & 11.8 & 0.997 \\
$\alpha \sim \mathcal{N}(0, 1^2)$ & 77.3 & 12.2 & 0.999 \\
\bottomrule
\end{tabular}
\end{table}

Results show minimal sensitivity: correlations $> 0.995$ across all variants.

\subsubsection*{A.4.2 Benchmark Subset Selection}

We test the impact of including/excluding specific benchmarks:

\begin{table}[h]
\centering
\caption{Ablation Study: Benchmark Contributions}
\begin{tabular}{lcc}
\toprule
\textbf{Benchmark Removed} & \textbf{$\rho$ w/ Arena ELO} & \textbf{$\Delta\rho$} \\
\midrule
None (Full model) & 0.89 & --- \\
HumanEval & 0.87 & -0.02 \\
LiveCodeBench & 0.83 & -0.06 \\
SciCode & 0.88 & -0.01 \\
Arena Coding Rank & 0.88 & -0.01 \\
Intelligence Index (aux) & 0.89 & 0.00 \\
\bottomrule
\end{tabular}
\end{table}

LiveCodeBench has the largest impact when removed, confirming its high loading ($\lambda = 0.96$).

\subsubsection*{A.4.3 Sample Size Requirements}

We investigate minimum sample size (number of models) needed for stable estimates:

\begin{figure}[h]
\centering
\includegraphics[width=0.6\textwidth]{fig_sample_size.pdf}
\caption{Posterior standard deviation vs. sample size. Estimates stabilize at N $\approx$ 100 models.}
\end{figure}

Results: With N $> 100$ models, posterior SDs are stable. Below N=50, uncertainty increases substantially.

\subsection*{A.5 Comparison with Factor Analysis}

BLF is related to classical Confirmatory Factor Analysis (CFA), but differs in several ways:

\begin{table}[h]
\centering
\caption{BLF vs. Classical Factor Analysis}
\begin{tabular}{lll}
\toprule
\textbf{Aspect} & \textbf{BLF} & \textbf{Classical CFA} \\
\midrule
Inference & Bayesian (MCMC) & Frequentist (ML/EM) \\
Missing data & Full Bayesian imputation & Listwise/FIML \\
Uncertainty & Full posterior & Asymptotic SE \\
Priors & Weakly informative & N/A \\
Interpretation & Composite scores & Latent constructs \\
Auxiliary variables & Explicit covariance & Requires separate model \\
\bottomrule
\end{tabular}
\end{table}

The Bayesian approach provides natural handling of missing data and full uncertainty quantification, which are critical for our application.

\subsection*{A.6 Computational Details}

\subsubsection*{A.6.1 Scaling Properties}

\begin{itemize}
    \item \textbf{Time complexity:} $O(N \cdot B \cdot S)$ where N=models, B=benchmarks, S=MCMC samples
    \item \textbf{Memory:} $O(N \cdot B + S \cdot (N + 3B))$ for data + posterior samples
    \item \textbf{Parallelization:} Chains run in parallel; 4 cores recommended
\end{itemize}

Typical runtime: 3-5 minutes for 250 models $\times$ 5 benchmarks on a modern CPU.

\subsubsection*{A.6.2 Implementation}

Available at: \texttt{https://github.com/yourusername/llm\_jury}

\begin{verbatim}
# Python example
from llm_jury.analysis.latent_factor import (
    CODING_BENCHMARKS,
    extract_benchmark_matrix,
    fit_latent_factor_model,
    summarize_latent_scores,
)

# Load data
df_scores, df_z, names, benchmarks = extract_benchmark_matrix(
    models, CODING_BENCHMARKS.get_configs()
)

# Fit model
idata = fit_latent_factor_model(
    z_obs, idx_model, idx_bench, n_models, n_benchmarks,
    draws=2000, tune=2000, chains=4
)

# Extract scores
df_result = summarize_latent_scores(idata, names, score_name='ccs')
\end{verbatim}

\subsection*{A.7 Limitations and Future Work}

\subsubsection*{A.7.1 Current Limitations}

\begin{enumerate}
    \item \textbf{Single latent factor:} Assumes a unidimensional quality construct. Extensions to multi-factor models (e.g., separate factors for code generation vs. debugging) are possible but increase complexity.
    
    \item \textbf{Linear relationships:} The model assumes linear relationships between latent quality and observed benchmarks. Nonlinear relationships could be captured with splines or neural networks, at the cost of interpretability.
    
    \item \textbf{Independence assumption:} Benchmark scores are assumed conditionally independent given $\theta_i$. In reality, some benchmarks may share task-specific variance (e.g., both require long-context reasoning).
    
    \item \textbf{Computational cost:} MCMC sampling is slower than frequentist methods. For real-time applications, scores must be precomputed.
\end{enumerate}

\subsubsection*{A.7.2 Future Extensions}

\begin{itemize}
    \item \textbf{Dynamic updating:} Online algorithms (e.g., Particle MCMC) for updating posteriors as new benchmark data arrives
    \item \textbf{Multi-factor models:} Separate latent factors for different skill dimensions
    \item \textbf{Hierarchical priors:} Group benchmarks by type (automated vs. human evaluation) with shared hyperpriors
    \item \textbf{Temporal dynamics:} Model drift in model quality over time (e.g., as training data improves)
\end{itemize}

\subsection*{References}

\begin{thebibliography}{99}

\bibitem{hoffman2014nuts}
Hoffman, M. D., \& Gelman, A. (2014). The No-U-Turn sampler: adaptively setting path lengths in Hamiltonian Monte Carlo. \textit{Journal of Machine Learning Research}, 15(1), 1593-1623.

\bibitem{gelman1992inference}
Gelman, A., \& Rubin, D. B. (1992). Inference from iterative simulation using multiple sequences. \textit{Statistical Science}, 7(4), 457-472.

\bibitem{betancourt2016diagnosing}
Betancourt, M. (2016). Diagnosing suboptimal cotangent disintegrations in Hamiltonian Monte Carlo. \textit{arXiv preprint arXiv:1604.00695}.

\bibitem{rubin1987multiple}
Rubin, D. B. (1987). \textit{Multiple imputation for nonresponse in surveys}. John Wiley \& Sons.

\bibitem{gelman2006data}
Gelman, A., \& Hill, J. (2006). \textit{Data analysis using regression and multilevel/hierarchical models}. Cambridge University Press.

\bibitem{embretson2000item}
Embretson, S. E., \& Reise, S. P. (2000). \textit{Item response theory for psychologists}. Lawrence Erlbaum Associates.

\end{thebibliography}
